% tps/tp5.tex : Compression : quantification
% Notebook associé : edge-ai-toolchain/notebooks/compression/05_quantization.py

\chapter{Compression : quantification}

\section*{Motivation}

Un modèle entraîné stocke ses poids en \texttt{float32}, 4 octets par paramètre. Les poids d'un réseau se concentrent autour de zéro, sur une plage étroite : 8 bits suffisent en général à les représenter avec une perte de précision faible. La quantification int8 divise la taille des poids par 4 et remplace les calculs flottants par des calculs entiers. L'ESP32-S3 dispose d'une unité de calcul flottant simple précision ; le gain en vitesse de l'int8 vient de la mémoire (4 fois moins d'octets à lire) et des instructions vectorielles 8 bits du processeur, exploitées par les bibliothèques d'inférence d'Espressif. Il se mesure au TP 7. Ce TP quantifie le MobileNetV2 du TP 4, d'abord après l'entraînement (PTQ, post-training quantization), puis pendant l'entraînement (QAT, quantization-aware training), et mesure ce que chacune coûte en précision. \\

\begin{imp}
Ce TP utilise le notebook \ToolchainLink{notebooks/compression/05_quantization.py}{\texttt{05\_quantization.py}}, qui part du modèle \texttt{models/cifar4\_mobilenetv2.pt} produit au TP 4. Les modèles int8 de PyTorch s'exécutent sur CPU (moteur \texttt{onednn} sous Windows) : les latences sont mesurées sur le CPU du PC, pas sur la carte. L'API \texttt{torch.ao.quantization} utilisée est en cours de remplacement par \texttt{torchao} ; pour la carte, la quantification finale passe par l'outil du framework embarqué (TP 7), sur le même principe.
Les valeurs citées dans les encadrés "À retenir" ont été mesurées sur un PC portable (GPU NVIDIA RTX 5060 Laptop, PyTorch 2.11) ; sur une autre machine ou sur CPU, elles varient de quelques dixièmes de point, et les durées dépendent du matériel.
\end{imp}

\section{Quantification d'une couche à la main}

La quantification affine \cite{jacob2018quantization} représente un réel $w$ par un entier $q$ sur 8 bits signés :
\[
q = \mathrm{clip}\left(\mathrm{round}\left(\frac{w}{s}\right) + z,\ -128,\ 127\right), \qquad \hat{w} = (q - z)\,s,
\]
où $s$ (\texttt{scale}) est le pas de quantification et $z$ (\texttt{zero\_point}) l'entier qui représente le réel 0. \\

\TQ{Pour des valeurs dans $[w_{min}, w_{max}]$, exprimer $s$ et $z$ de façon à utiliser les 256 niveaux. Quelle est l'erreur maximale de reconstruction $|\hat{w} - w|$ ?}

\begin{figure}[H]
\centering
\resizebox{0.78\linewidth}{!}{% diagrams/tp5/quantization.tex : correspondance entre l'axe réel et la grille int8
\begin{tikzpicture}[font=\sffamily\small, >=Stealth]
% Axe réel
\draw[->, thick] (-0.5,1.8) -- (12.5,1.8) node[right] {réel $w$};
\foreach \x/\t in {1/w_{min}, 4.2/0, 11/w_{max}} {
    \draw[thick] (\x,1.65) -- (\x,1.95);
    \node[above] at (\x,1.95) {$\t$};
}
% Grille entière
\draw[->, thick] (-0.5,0) -- (12.5,0) node[right] {entier $q$};
\foreach \k in {0,...,20} \draw (1+\k*0.5,-0.08) -- (1+\k*0.5,0.08);
\node[below] at (1,-0.1) {$-128$};
\node[below] at (4.2,-0.1) {$z$};
\node[below] at (11,-0.1) {$127$};
% Correspondances exactes
\foreach \x in {1, 4.2, 11} \draw[->, ECE, thick] (\x,1.62) -- (\x,0.18);
% Une valeur quelconque est arrondie au niveau le plus proche
\draw[->, gamboge, thick] (7.3,1.62) -- (7.5,0.18);
\node[gamboge!80!black, font=\sffamily\scriptsize, align=center] at (7.4,2.3) {valeur quelconque :\\arrondie au niveau le plus proche};
\draw[gamboge, thick] (7.3,1.65) -- (7.3,1.95);
% Pas de quantification
\draw[decorate, decoration={brace, amplitude=3pt, mirror}, black!60] (2,-0.15) -- (2.5,-0.15)
    node[midway, below=3pt, font=\sffamily\scriptsize, black!70] {pas $s$};
\node[font=\sffamily\scriptsize, black!70, align=left, anchor=west] at (-0.4,-1.0)
    {256 niveaux couvrent $[w_{min}, w_{max}]$ : $s = (w_{max} - w_{min}) / 255$ ; le réel 0 tombe exactement sur l'entier $z$.};
\end{tikzpicture}
}
\caption{Quantification affine : la plage réelle des poids est découpée en 256 niveaux de pas $s$ ; le réel 0 correspond à l'entier $z$.}
\label{fig:tp5-quant}
\end{figure}

\EQ{Compléter \texttt{quantize\_per\_tensor\_numpy()} et l'appliquer aux poids de la couche \texttt{features[6].conv[2]}. Relever la plage des poids, $s$, $z$, l'erreur maximale et l'erreur moyenne. Vérifier que le résultat est identique à \texttt{torch.quantize\_per\_tensor()}.}

\EQ{Tracer l'histogramme des poids de cette couche. Expliquer avec cet histogramme pourquoi 256 niveaux suffisent.}

\EQ{Le notebook calcule aussi une quantification par canal (un $s$ par canal de sortie, symétrique). Comparer l'erreur moyenne à celle de la quantification par tenseur. Pourquoi est-elle plus faible ?}

\section{Quantification après entraînement (PTQ)}

La PTQ se déroule en trois étapes : \texttt{fuse\_model()} fusionne chaque convolution avec la BatchNorm et l'activation qui la suivent ; des observateurs enregistrent la plage des activations sur un jeu de calibration ; \texttt{convert()} remplace les couches par leurs versions int8. \\

\TQ{Pourquoi peut-on fusionner une convolution et la BatchNorm qui la suit en une seule convolution, une fois l'entraînement terminé ? Quel gain en tire l'inférence ?}

\TQ{Les poids sont connus après l'entraînement ; les activations dépendent de l'image d'entrée. Pourquoi la PTQ a-t-elle besoin d'un jeu de calibration ? D'où doivent venir ces images ?}

\EQ{Quantifier le modèle avec 128 images de calibration. Relever taille, précision de test et latence CPU du modèle \texttt{float32} et du modèle int8. Calculer le rapport de taille. Pourquoi n'est-il pas exactement de 4 ?}

\EQ{Le notebook refait la calibration avec 8 images, avec trois tirages différents de 128 images, puis avec 128 images d'une seule classe. Relever la précision de test dans chaque cas. Le jeu de test compte 400 images : de combien de points varie la précision quand une seule image change de classe ? Les écarts mesurés entre jeux de calibration sont-ils significatifs ?}

\section{Entraînement avec quantification simulée (QAT)}

La QAT insère dans le réseau des opérations "fake quant" qui arrondissent poids et activations comme le fera l'int8, tout en calculant en flottant. Le réseau, réentraîné quelques epochs, s'adapte à l'erreur de quantification. \\

\EQ{Réentraîner le modèle 3 epochs en QAT (taux d'apprentissage $10^{-5}$, BatchNorm figées), puis le convertir en int8. Relever précision, taille et latence, et compléter le tableau de synthèse du notebook.}

\TQ{Sur ce modèle, la QAT récupère une partie de la précision perdue par la PTQ. Combien de points, et pour quel coût en temps d'entraînement ? Dans quels cas la QAT devient-elle nécessaire ?}

\begin{imp}
\textbf{À retenir :} la PTQ int8 divise ici la taille du modèle par 3.5 (9.16~Mo à 2.65~Mo) pour 1.7 point de précision perdu (88.5~\% à 86.8~\%), sans réentraînement ; la QAT en récupère 0.4 (87.2~\%). Sur 400 images de test, une image vaut 0.25 point : un écart d'un point entre deux réglages reste à la limite de la mesure. Commencer par la PTQ ; ne passer à la QAT que si la perte mesurée est trop grande pour l'application.
\end{imp}
